% Section 9.3, from lectures/L09/README.md: what you should be able to do afterwards, the
% questions to test yourself, and the next lecture.

\section{Review}
\label{c9:sec:review}

\needspace{6\baselineskip}
\subsection*{What you should be able to do now}
After this chapter, you should be able to:
\begin{itemize}
  \item Say what state a sleeping task is in, and what a spinning one is in.
  \item Describe the lost wakeup, and say which line of \code{wait_event_interruptible} prevents it.
  \item Explain why the condition passed to a wait macro is evaluated more than once.
  \item Implement a blocking \code{read()} that handles \code{O_NONBLOCK} and signals correctly.
  \item Implement \code{poll}, and say what your driver must do for \code{select} to work on it.
  \item Choose between \code{udelay} and \code{msleep} from the context and the duration.
  \item Say what a jiffy is on a given kernel, and why quoting a duration in jiffies is a mistake.
  \item Choose between a \code{timer_list} and an \code{hrtimer}, and say what resolution each gives
    you.
\end{itemize}

\needspace{6\baselineskip}
\subsection*{Questions to test yourself}
\begin{enumerate}
  \item Why does \code{wait_event_interruptible} take a condition rather than a variable to wait on?
  \item Your reader wakes, finds no data, and returns 0 to userspace. What did the caller conclude,
    and was it true?
  \item A user presses Ctrl-C while your \code{read()} is blocked and nothing happens. What did you
    fail to return?
  \item What is \code{HZ} on the kernel you built, and how many jiffies is a 3 ms delay?
  \item \code{msleep(1)} in a loop of 1000 takes noticeably more than one second. Why?
  \item Why can \code{udelay(5000)} be a bug even in code that is allowed to busy-wait?
  \item Your driver implements \code{read} but not \code{poll}. What happens when a program calls
    \code{select} on it?
\end{enumerate}

\needspace{6\baselineskip}
\subsection*{Where to look}
\begin{itemize}
  \item Chapter~\ref{c8:ch} produced the events this chapter waits for. The two labs are one driver.
\end{itemize}

\needspace{6\baselineskip}
\subsection*{Looking ahead}
Chapter~\ref{c10:ch} is about:
\begin{itemize}
  \item How a driver finds its hardware without an address in its source.
  \item Bus, device and driver, and the match that puts the three together.
  \item Device tree syntax, properly, and the \code{ranges} you already translated by hand in
    Chapter~\ref{c6:ch}.
  \item \code{probe}, and the moment your driver stops being a module that pokes at an address.
\end{itemize}
